\documentclass[10pt]{article}
\usepackage[T1]{fontenc}
\usepackage[utf8]{inputenc}
\usepackage[french]{babel}
\usepackage{amsmath,amssymb,mathrsfs,mathpazo}
\usepackage{geometry}
\geometry{a4paper,margin=2cm}
\usepackage{graphicx,array,tabularx,tabularray,colortbl,multirow,multicol,booktabs,enumitem,pifont,ulem,fancybox,diagbox,cancel}
\usepackage{pstricks,pstricks-add,pst-plot,pst-tree,pst-node,pst-text,pst-func,pst-math,pst-bspline,pst-3dplot,pst-eucl,pst-solides3d}
\let\pscalculate\undefined
\usepackage{pst-bezier}
\usepackage{tikz,pgfplots,tkz-tab}
\usetikzlibrary{arrows,arrows.meta,patterns,calc,positioning,trees,decorations.markings,shapes.geometric,shapes.misc,shapes.arrows,intersections,decorations.pathmorphing,decorations.pathreplacing,backgrounds,shadows,babel}
\pgfplotsset{compat=1.18}
\usepgfplotslibrary{fillbetween}
\usepackage[np]{numprint}
\usepackage{listings,makecell,pas-tableur,twoopt}
\let\barre\undefined
\usepackage[pstricks]{tablvar}
\definecolor{darkgreen}{RGB}{0,100,0}
\providecommand{\red}{\color{red}}
\providecommand{\blue}{\color{blue}}
\providecommand{\green}{\color{green}}
\providecommand{\black}{\color{black}}
\providecommand{\white}{\color{white}}
\providecommand{\gray}{\color{gray}}
\providecommand{\cyan}{\color{cyan}}
\providecommand{\magenta}{\color{magenta}}
\providecommand{\yellow}{\color{yellow}}
\providecommand{\violet}{\color{violet}}
\providecommand{\orange}{\color{orange}}
\providecommand{\euro}{\texteuro}
\providecommand{\EUR}{\texteuro}
\providecommand{\ieme}{\textsuperscript{e}}
\providecommand{\degres}{\ensuremath{^\circ}}
\providecommand{\degree}{\ensuremath{^\circ}}
\providecommand{\vv}[1]{\overrightarrow{#1}}
\providecommand{\Coord}[1]{#1}
\providecommand{\decosix}{}
\providecommand{\starredbullet}{\ensuremath{\bullet}}
\pagestyle{empty}
\setlength{\parindent}{0pt}
\usepackage{adjustbox}\DeclareUnicodeCharacter{2236}{:}\DeclareUnicodeCharacter{2190}{\ensuremath{\leftarrow}}
\providecommand{\R}{\mathbb{R}}
\providecommand{\N}{\mathbb{N}}
\providecommand{\Z}{\mathbb{Z}}
\providecommand{\Q}{\mathbb{Q}}
\providecommand{\D}{\mathbb{D}}
\providecommand{\C}{\mathbb{C}}
\providecommand{\e}{{\mathrm{e}}}
\providecommand{\vect}[1]{\overrightarrow{#1}}
\providecommand{\vectt}[1]{\overrightarrow{\mathrm{#1}}}
\providecommand{\vvt}[1]{\overrightarrow{\mathrm{#1}}}
\providecommand{\np}[2][]{{#2~#1}}
\providecommand{\vv}[1]{{\overrightarrow{#1}}}
\renewcommand{\d}{{\mathrm{d}}}
\providecommand{\degres}{{^\circ}}
\providecommand{\degre}{{^\circ}}
\providecommand{\blue}{\color{blue}}
\providecommand{\red}{\color{red}}
\providecommand{\enskip}{\quad}
\providecommand{\llbracket}{\lbrack\!\lbrack}
\providecommand{\rrbracket}{\rbrack\!\rbrack}
\providecommand{\euro}{{\text{€}}}
\providecommand{\no}{{\text{n°}}}
\providecommand{\interval}[2]{{[#1;#2]}}
\providecommand{\intervFO}[2]{{[#1;#2[}}
\providecommand{\dotfill}{\ldots}

\providecommand{\ds}{\displaystyle}
\providecommand{\textding}[1]{#1}
\providecommand{\Oij}{\(\mathrm{O};\vec{\imath},\vec{\jmath}\)}
\providecommand{\Oijk}{\(\mathrm{O};\vec{\imath},\vec{\jmath},\vec{k}\)}
\providecommand{\Ouv}{\(\mathrm{O};\vec{u},\vec{v}\)}
\providecommand{\intervFF}[2]{\ensuremath{[#1;#2]}}
\renewcommand{\euro}{\ifmmode\text{\texteuro}\else\texteuro\fi}
\DeclareUnicodeCharacter{00B0}{\ifmmode^{\circ}\else\textdegree\fi}
\DeclareUnicodeCharacter{20AC}{\ifmmode\text{\texteuro}\else\texteuro\fi}
% Énoncés : Éducation nationale. Transcription et corrigés : APMEP, Denis Vergès et contributeurs. Découpage et classement : CoopMaths, Sébastien Lozano, Éric Elter et contributeurs. Adaptation de la présentation : Maths Explorer.
% CC BY-SA : https://coopmaths.fr/www/about/
% https://coopmaths.fr/alea/static/bac/2024/tex/bac_2024_06_sujet2_polynesie_4.tex
\begin{document}
\noindent\textit{APMEP / CoopMaths. CC BY-SA. Présentation adaptée par Maths Explorer.}\par\medskip
\medskip

Une commune décide de remplacer le traditionnel feu d'artifice du 14 juillet par un
spectacle de drones lumineux.

\smallskip

Pour le pilotage des drones, l'espace est muni d'un repère orthonormé \Oijk{} dont l'unité est la centaine de mètres.

\smallskip

La position de chaque drone est modélisée par un point et chaque drone est envoyé d'un point de départ D de coordonnées $(2~;~5~;~1)$.

\smallskip

On souhaite former avec des drones des figures en les positionnant dans un même plan $\mathcal{P}$.

\smallskip

Trois drones sont positionnés aux points A$(-1~;~-1~;~17)$, B$(4~;~-2~;~4)$ et C$(1~;~-3~;~7)$.

\begin{enumerate}
	\item Justifier que les points A, B et C ne sont pas alignés.
\end{enumerate}

Dans la suite, on note $\mathcal{P}$ le plan (ABC) et on considère le vecteur $\vect{n} \begin{pmatrix}2\\-3\\1\end{pmatrix}$.

\begin{enumerate}[resume]
	\item 
	\begin{enumerate}
		\item Justifier que $\vect{n}$ est normal au plan (ABC).
		\item Démontrer qu'une équation cartésienne du plan $\mathcal{P}$ est $2x - 3y + z - 18 = 0$.
	\end{enumerate}
	\item Le pilote des drones décide d'envoyer un quatrième drone en prenant comme trajectoire la droite $d$ dont une représentation paramétrique est donnée par 

\[ d~:~\left\{\begin{array}{l c l}
x&=&3t+2\\y&=&\phantom{3}t+5\\z&=&4t+1\end{array}\right.\text{, avec } t \in \R. \]
	\begin{enumerate}
		\item Déterminer un vecteur directeur de la droite $d$.
		\item Afin que ce nouveau drone soit également placé dans le plan $\mathcal{P}$, déterminer par le calcul les coordonnées du point E, intersection de la droite $d$ avec le plan $\mathcal{P}$.
	\end{enumerate}
	\item Le pilote des drones décide d'envoyer un cinquième drone le long de la droite $\Delta$ qui passe par le point $D$ et qui est perpendiculaire au plan $\mathcal{P}$. Ce cinquième drone est placé lui aussi dans le plan $\mathcal{P}$, soit à l'intersection entre la droite $\Delta$ et le plan $\mathcal{P}$. On admet que le point F$(6~;~-1~;~3)$ correspond à cet emplacement.
	
Démontrer que la distance entre le point de départ D et le plan $\mathcal{P}$ vaut $2\sqrt{14}$ centaines de mètres.
	\item L'organisatrice du spectacle demande au pilote d'envoyer un nouveau drone dans le plan (peu importe sa position dans le plan), toujours à partir du point D.
	
Sachant qu'il reste $40$~secondes avant le début du spectacle et que le drone vole en trajectoire rectiligne à 18,6~$\text{m.s}^{-1}$, le nouveau drone peut-il arriver à temps ?
\end{enumerate}

\end{document}
